\documentclass[11pt,letterpaper]{article}
\usepackage[utf8]{inputenc}
\usepackage[T1]{fontenc}
\usepackage[spanish,es-nodecimaldot,es-tabla]{babel}
\usepackage{lmodern}
\usepackage[margin=2.4cm]{geometry}
\usepackage{amsmath,amssymb,amsthm}
\usepackage{bm}
\usepackage{graphicx}
\usepackage{booktabs}
\usepackage{caption}
\usepackage{enumitem}
\usepackage[hidelinks]{hyperref}

\newcommand{\avg}[1]{\left\langle #1 \right\rangle}
\newcommand{\kB}{k_{B}}
\newcommand{\dd}{\mathrm{d}}
\newtheorem*{lema}{Lema}
% Si falta una figura o tabla generada por el código, se muestra un recuadro en su lugar
\newcommand{\figura}[2][]{\IfFileExists{#2}{\includegraphics[#1]{#2}}%
  {\fbox{\parbox{0.6\textwidth}{\centering\vspace{2em}Figura pendiente:\\\texttt{\detokenize{#2}}\vspace{2em}}}}}
\newcommand{\tablaopcional}[1]{\IfFileExists{#1}{\input{#1}}%
  {\begin{center}\fbox{Tabla pendiente: \texttt{\detokenize{#1}}}\end{center}}}

\title{\textbf{Avance de la semana 12}\\[4pt]
\large Método de Monte Carlo en el ensamble microcanónico para sistemas discretos:\\
reproducción de Ray y Frelechoz, \textit{Phys.\ Rev.\ E} \textbf{53}, 3402 (1996)}
\author{Andrés Felipe Contreras \and Melany Daza Valencia\\[4pt]
\small Física Estadística e Introducción a la Teoría de la Información --- Universidad EAFIT, 2026-2\\
\small Profesores: Daniela Angulo y Luis Eduardo Olmos Sánchez}
\date{}

\begin{document}
\maketitle

% \begin{abstract}
% \noindent
% Este avance contiene (i) un resumen del artículo de Ray y Frelechoz~\cite{ray96},
% (ii) el planteamiento de su resultado central, (iii) la derivación analítica completa
% de todas las fórmulas del artículo, con los pasos intermedios que el texto omite, y
% (iv) un código preliminar que implementa los métodos de Monte Carlo canónico y
% microcanónico para el modelo de Ising bidimensional, con una prueba de verificación
% interna contra la enumeración exacta de una red $4\times4$ y una primera comparación
% con las Tablas~I--II y las Figuras~1--3 del artículo para la red $30\times30$.
% \end{abstract}

%=====================================================================
\section{Resumen del artículo}
%=====================================================================

\paragraph{Problema y contexto.}
Desde el trabajo de Metropolis \emph{et al.}~\cite{metropolis53}, casi todas las
simulaciones de Monte Carlo (MC) se hacen en el ensamble canónico, finjando la
temperatura $T$ y promediando la energía y sus fluctuaciones. El ensamble
microcanónico, que describe un sistema aislado y es en ese sentido el más
fundamental, se ha utilizado de forma más limitada porque en trabajo analítico es más incómodo.
Para sistemas con potenciales continuos el ensamble microcanónico se genera con
dinámica molecular, pero en un sistema discreto (espines, gas de red, aleaciones)
no hay ecuaciones de movimiento. La alternativa usada hasta entonces era el
algoritmo del ``demonio'' de Creutz~\cite{creutz83}, muy eficiente, pero que según
los autores no es un método de MC con fundamento riguroso en la mecánica
estadística y no permite calcular directamente funciones respuesta como el calor
específico y la susceptibilidad mediante fórmulas de fluctuación. En un trabajo
previo, Ray~\cite{ray91} había propuesto un MC microcanónico para potenciales
continuos, el artículo extiende ese método a sistemas discretos.

\paragraph{Idea central.}
A cada espín se le asocia una partícula ficticia de masa muy grande con momento
$\bm p_i\in\mathbb R^{d}$, de modo que el hamiltoniano total es
$H(\bm p,s)=K(\bm p)+U(s)$, con $U$ la energía de Ising
\begin{equation}
U(s)=-J\sum_{\langle ij\rangle}S_iS_j-B\sum_i S_i ,\qquad S_i=\pm1 .
\label{eq:U}
\end{equation}
La densidad microcanónica es uniforme sobre la superficie $H=E$,
$W_m(\bm p,s)=D\,\delta(E-H)$. Al integrar sobre los $dN$ momentos se obtiene una
densidad de probabilidad sobre las configuraciones de espín:
\begin{equation}
W_E(s)=D'\,\bigl(E-U(s)\bigr)^{dN/2-1}\,\Theta\bigl(E-U(s)\bigr),
\label{eq:WE}
\end{equation}
donde $\Theta$ es una función escalón y $D'$ es una constante. Esta densidad se muestrea con el algoritmo de Metropolis, se propone invertir
un espín recorriendo la red en orden y se acepta con probabilidad
$\min[1,\,W_E(s')/W_E(s)]$. La energía cinética promedio define la temperatura,
$dN\kB T/2=\avg{K}=E-\avg{U}$, y la derivada de $T$ respecto a $E$ da fórmulas de
fluctuación propias del ensamble microcanónico. Para el calor específico,
\begin{equation}
\frac{dN\kB}{2C_B}=\frac{dN}{2}-\Bigl(\frac{dN}{2}-1\Bigr)\avg{K}\avg{\frac1K},
\label{eq:CB}
\end{equation}
donde $C_B=dN\kB/2+C$ incluye la contribución ``gas ideal'' de las partículas
ficticias y $C$ es el calor específico configuracional. A diferencia del caso
canónico, $C=(\avg{U^2}-\avg{U}^2)/\kB T^2$, la expresión microcanónica no puede
escribirse sólo en términos de $U$. También se dan expresiones para
la susceptibilidad a energía constante $\chi_E$ y las relaciones que conectan $\chi_E$ con la
susceptibilidad isotérmica $\chi_T$ del ensamble canónico a partir de cantidades evaluadas en el ensamble microcanónico.

\paragraph{Simulaciones.}
Se estudia el modelo de Ising bidimensional en una red cuadrada de
$N = 30\times30$ espines $S_i=\pm1$, con condiciones de frontera periódicas y
sin campo externo ($B=0$), de modo que la energía potencial es
$U(s) = -J\sum_{\langle ij\rangle} S_i S_j$. Se trabaja en unidades
$J = \kB = 1$, por lo que $T$ se mide en unidades de $J/\kB$ y el calor
específico por espín es $c = C/N$. En el ensamble canónico se eligen nueve temperaturas
entre $1.5$ y $3.0$; se equilibra durante al menos $5\times10^6$ barridos y luego se
hacen cinco corridas independientes de $10^6$ barridos, cuya desviación estándar
se reporta como incertidumbre. Con el $\avg{U}$ canónico se calcula
$E=N\kB T+\avg U$ (Ec.~\eqref{eq:E-can} con $d=2$), y esas nueve energías son la
entrada de las simulaciones microcanónicas, hechas con el mismo protocolo.

\paragraph{Resultados.}
\begin{itemize}[leftmargin=1.2em,itemsep=1pt]
\item La curva $E(T)$ (Fig.~1) es prácticamente idéntica en ambos ensambles, como
se espera de la equivalencia de ensambles (Lebowitz, Percus y Verlet~\cite{lpv67}).
\item La magnetización media $\avg{|M|}/|M_0|$ (Fig.~4) también coincide.
\item El calor específico configuracional $c=C/N\kB$ (Figs.~2--3) coincide
en todo el intervalo excepto cerca del máximo: en $T\approx2.30$ el valor
microcanónico es $c=2.1707\pm0.03$ y el canónico $c=1.8712\pm0.004$, una diferencia
del 16~\%, muy superior a la diferencia de orden $1/N$ (alrededor de 0.1~\%) esperada por tamaño
finito.
\item En una red $60\times60$ el máximo se desplaza a $T\approx2.28$ (hacia el valor
de Onsager $T_c=2.269$) y los calores específicos son $2.464\pm0.02$ (micro) y
$2.226\pm0.02$ (canónico): la diferencia baja a 10~\%.
\end{itemize}

\paragraph{Conclusiones.}
El MC microcanónico tiene el mismo costo computacional que el canónico, permite
calcular cualquier propiedad termodinámica mediante fórmulas de fluctuación, y es
la única forma de generar el ensamble microcanónico en un sistema discreto. La
gran diferencia entre ensambles en el calor específico es un precursor de la
transición de fase del límite termodinámico, y se sugiere que podría usarse
para localizar transiciones de fase en otros sistemas.

% \paragraph{Lectura crítica preliminar.}
% El ``ensamble microcanónico'' del artículo es el del sistema compuesto espines $+$
% gas ideal ficticio de $dN$ grados de libertad. Desde el punto de vista de los espines,
% el gas actúa como un baño térmico \emph{finito}, con capacidad calorífica
% $dN\kB/2$ comparable a la de los espines. Por eso las fluctuaciones de $U$ quedan
% suprimidas respecto al caso canónico (Sec.~\ref{sec:equiv}) y la diferencia entre
% ensambles se amplifica justo donde $C$ es grande. El número de grados de libertad
% ficticios, $n=dN/2$, es en realidad un parámetro libre del método: $n=1$ se parece al
% demonio de Creutz y $n\to\infty$ recupera el canónico. El artículo no explora esta
% dependencia; la incluimos como extensión natural del proyecto.  

\paragraph{Resultado que se va a reproducir.}
Las Tablas~I y II y las Figuras~1--4: la curva $E(T)$ y el calor específico $c(T)$
en ambos ensambles para la red $30\times30$, en particular la diferencia del
16~\% en el máximo ($T\approx2.30$).

%=====================================================================
\section{Derivación analítica}
%=====================================================================

Partiendo únicamente de los postulados de los ensambles canónico,
$W_c(\bm p,s)\propto e^{-H/\kB T}$, y microcanónico, $W_m(\bm p,s)\propto\delta(E-H)$,
con $H=\sum_{a=1}^{dN}p_a^2/2m+U(s)$ y $U$ dado por~\eqref{eq:U}, se debe demostrar:
\begin{enumerate}[label=(\alph*),itemsep=1pt]
\item Ensamble canónico: la densidad de configuraciones $W_T(s)\propto e^{-U/\kB T}$,
la equipartición $\avg E=dN\kB T/2+\avg U$ y las fórmulas de fluctuación para $C$
y $\chi_T$ [Ecs.~(2), (4), (6), (7) del artículo].
\item Ensamble microcanónico: la densidad~\eqref{eq:WE} [Ec.~(9)], que la regla de
Metropolis con esa densidad la deja invariante, y que la definición
$dN\kB T/2=\avg K$ [Ec.~(11)] es la temperatura termodinámica.
\item La fórmula de fluctuación microcanónica del calor
específico~\eqref{eq:CB} a utiliazar en las simulaciones [Ec.~(12)].
\item La susceptibilidad a energía constante $\chi_E$ [Ec.~(13)] y las relaciones
entre $\chi_E$ y $\chi_T$ en cada ensamble [Ecs.~(14) y (15)].
\item Que en el límite termodinámico ambos ensambles dan el mismo $c$, y por qué en
un sistema finito difieren.
\end{enumerate}

%=====================================================================
\section{Derivación analítica}
%=====================================================================

$N$ es el número de espines, $d$ la dimensión espacial, $f=dN$ el número de grados de
libertad de momento y $n=f/2=dN/2$. Escribimos $K=E-U(s)$ para la energía cinética
en el ensamble microcanónico. Como $U$ depende de $B$ linealmente\eqref{eq:U}:
\begin{equation}
\frac{\partial U}{\partial B}=-\sum_iS_i\equiv-M,\qquad \frac{\partial M}{\partial B}=0 ,
\label{eq:dUdB}
\end{equation}

donde $M$ es la magnetización total. 

%---------------------------------------------------------------------
\subsection{Ensamble canónico}
%---------------------------------------------------------------------
Como $H$ es separable, la integral sobre momentos se factoriza con $K=\sum_a p_a^2/2m$:
\begin{equation}
\int\dd^{f}p\;e^{-\beta\sum_a p_a^2/2m}=\prod_{a=1}^{f}\int_{-\infty}^{\infty}\dd p_a\,
e^{-\beta p_a^2/2m}=\Bigl(\frac{2\pi m}{\beta}\Bigr)^{f/2},\qquad\beta=\frac1{\kB T},
\end{equation}
que no depende de $s$. Por lo tanto se absorbe en la constante de normalización $A'$, y la densidad marginal de configuraciones es
$W_T(s)=A'e^{-\beta U(s)}$ [Ec.~(2)], con la función de partición
$Z=\sum_s e^{-\beta U(s)}$. La energía
cinética media es
$\avg K=-\partial_\beta\ln(2\pi m/\beta)^{f/2}=f/2\beta$, de donde
\begin{equation}
\avg E=\frac{dN\kB T}{2}+\avg U .
\label{eq:E-can}
\end{equation}
\textbf{Calor específico.} A $B$ fijo, $T$ sólo entra por $\beta=1/\kB T$. De
$\partial_\beta Z=-\sum_sUe^{-\beta U}=-Z\avg U$ se obtiene $\avg U=-\partial_\beta\ln Z$, y
derivando $\avg U=\sum_sUe^{-\beta U}/Z$ con la regla del cociente, se obtiene la expresión
$\partial_\beta\avg U=-\avg{U^2}+\avg U^2$. Con $\partial_T=(\dd\beta/\dd T)\partial_\beta=-\kB\beta^2\partial_\beta$:
\begin{equation}
C\equiv\Bigl(\frac{\partial\avg U}{\partial T}\Bigr)_B
=\kB\beta^2\bigl(\avg{U^2}-\avg U^2\bigr)
=\frac{\avg{U^2}-\avg U^2}{\kB T^2}.
\label{eq:C-can}
\end{equation}
Es la capacidad calorífica configuracional (la parte cinética aporta $dN\kB/2$ por
separado). En el código, con $J=\kB=1$, $c=\mathrm{Var}(U)/NT^2$.

\vspace{0.5cm}
\noindent
\textbf{Susceptibilidad.} Con~\eqref{eq:dUdB}, $\avg M=\beta^{-1}\partial_B\ln Z$, y
derivando el cociente $\sum_sMe^{-\beta U}/Z$:
\begin{equation}
\chi_T\equiv\Bigl(\frac{\partial\avg M}{\partial B}\Bigr)_T=\frac{\avg{M^2}-\avg M^2}{\kB T}.
\label{eq:chiT}
\end{equation}
Necesitaremos además dos derivadas mixtas para las relaciones entre $\chi_E$ y $\chi_T$, que salen del mismo procedimiento:
\begin{align}
\Bigl(\frac{\partial\avg M}{\partial T}\Bigr)_B&=\frac{\avg{MU}-\avg M\avg U}{\kB T^2},
\label{eq:dMdT}\\
\Bigl(\frac{\partial\avg U}{\partial B}\Bigr)_T&=\Bigl\langle\frac{\partial U}{\partial B}\Bigr\rangle
-\beta\Bigl(\Bigl\langle U\frac{\partial U}{\partial B}\Bigr\rangle-\avg U\Bigl\langle\frac{\partial U}{\partial B}\Bigr\rangle\Bigr)
=-\avg M+\frac{\avg{MU}-\avg M\avg U}{\kB T}.
\label{eq:dUdB-T}
\end{align}

%---------------------------------------------------------------------
\subsection{Ensamble microcanónico: densidad de configuraciones}
%---------------------------------------------------------------------
Integramos $W_m=D\,\delta(E-K-U)$ sobre los $f$ momentos usando coordenadas
hiperesféricas, $\dd^fp=S_f\,p^{f-1}\dd p$ con $S_f=2\pi^{f/2}/\Gamma(f/2)$ el área de
la esfera unitaria en $\mathbb R^f$ y $p=|\bm p|$. Para $E>U$ la delta tiene una
única raíz positiva $p_*=\sqrt{2m(E-U)}$ y
$\delta\bigl(E-U-p^2/2m\bigr)=(m/p_*)\,\delta(p-p_*)$. Entonces
\begin{equation}
\int\dd^fp\;\delta\Bigl(E-U-\frac{p^2}{2m}\Bigr)=S_f\,m\,p_*^{f-2}
=S_f\,m\,(2m)^{f/2-1}\,(E-U)^{f/2-1}\,\Theta(E-U),
\end{equation}
donde el escalón aparece porque para $E<U$ la delta no tiene raíz ($K\ge0$). Absorbiendo
las constantes en $D'$ se obtiene~\eqref{eq:WE}:
$W_E(s)=D'\,K^{\,n-1}\Theta(K)$ con $K=E-U(s)$. Los promedios microcanónicos son
$\avg A=\sum_sW_E(s)A(s)/\sum_sW_E(s)$.

\paragraph{Regla de aceptación.} La probabilidad de transición
$P(s\to s')=g(s\to s')\,a(s\to s')$, con probabilidad de la propuesta $g$ (invertir un espín)
y probabilidad de aceptarla $a=\min[1,W_E(s')/W_E(s)]$. Ya que la propuesta es simétrica, $g(s\to s') = g(s'\to s)$, (invertir el mismo espín te devuelve a $s$) se cumple el balance detallado,
$W_E(s)P(s\to s')=W_E(s')P(s'\to s)=g\,\min[W_E(s),\,W_E(s')]$, de modo que $W_E$ es
estacionaria. Recorrer los espines en orden compone pasos que individualmente dejan
$W_E$ invariante, así que la composición también la deja invariante. Explícitamente,
si $\Delta U=U(s')-U(s)$,
\begin{equation}
a(s\to s')=\min\Bigl[1,\Bigl(1-\frac{\Delta U}{K}\Bigr)^{n-1}\Theta(K-\Delta U)\Bigr]
\;\xrightarrow{\;n\to\infty\;}\;\min\bigl[1,e^{-\Delta U/\kB T}\bigr],
\label{eq:acc}
\end{equation}
con $\kB T=K/(n-1)$ fijo, lo que muestra que el algoritmo se reduce al Metropolis
canónico cuando el ``baño'' ficticio es infinito.

%---------------------------------------------------------------------
\subsection{Temperatura microcanónica}
%---------------------------------------------------------------------
Definimos las sumas
\begin{equation}
Z_k(E,B)\equiv\sum_s\bigl(E-U(s)\bigr)^k\,\Theta\bigl(E-U(s)\bigr),
\end{equation}
de modo que la normalización de $W_E$ es $Z_{n-1}$ y, por ejemplo,
$\avg K=Z_n/Z_{n-1}$ y $\avg{1/K}=Z_{n-2}/Z_{n-1}$.

\begin{lema}
Para $k\ge1$, $\partial_EZ_k=k\,Z_{k-1}$ y $\partial_BZ_k=k\sum_sM(s)(E-U)^{k-1}\Theta(E-U)$.
\end{lema}
\noindent\emph{Demostración.} Derivando término a término,
$\partial_E[(E-U)^k\Theta]=k(E-U)^{k-1}\Theta+(E-U)^k\delta(E-U)$, y el segundo término se
anula porque $(E-U)^k=0$ donde la delta tiene soporte. Para $B$ se usa
$\partial_B(E-U)=M$ por~\eqref{eq:dUdB}. 

El volumen de fase encerrado por la superficie $H=E$ es
$\Phi(E)=\int\dd^fp\sum_s\Theta(E-H)=V_f(2m)^{f/2}Z_n$, con $V_f$ el volumen de la bola
unitaria en $\mathbb R^f$. Con la entropía de Gibbs $S=\kB\ln\Phi$ y el lema,
\begin{equation}
\frac1T=\Bigl(\frac{\partial S}{\partial E}\Bigr)_B=\kB\frac{nZ_{n-1}}{Z_n}
\quad\Longrightarrow\quad
\frac{dN}{2}\kB T=\frac{Z_n}{Z_{n-1}}=\avg K ,
\label{eq:T}
\end{equation}
que es exactamente la Ec.~(11) del artículo. Si en cambio se usa la entropía de
superficie $S=\kB\ln\Omega$ con $\Omega=\partial_E\Phi\propto Z_{n-1}$, se obtiene
$1/\kB T=(n-1)\avg{1/K}$; ambas definiciones difieren en $O(1/N)$.

%---------------------------------------------------------------------
\subsection{Resultado central: calor específico microcanónico}
%---------------------------------------------------------------------
Derivamos~\eqref{eq:T} respecto a $E$ con el lema:
\begin{equation}
\frac{\partial\avg K}{\partial E}
=\frac{\partial_EZ_n}{Z_{n-1}}-\frac{Z_n\,\partial_EZ_{n-1}}{Z_{n-1}^2}
=\frac{nZ_{n-1}}{Z_{n-1}}-(n-1)\frac{Z_n}{Z_{n-1}}\frac{Z_{n-2}}{Z_{n-1}}
=n-(n-1)\avg K\avg{\frac1K}.
\end{equation}
Como $\avg K=n\kB T$, la capacidad calorífica total a $B$ constante,
$1/C_B=(\partial T/\partial E)_B$, satisface
\begin{equation}
\boxed{\;\frac{n\kB}{C_B}=n-(n-1)\avg K\avg{\frac1K}\;}\qquad
\Bigl(n=\frac{dN}{2}\Bigr),
\label{eq:CB-der}
\end{equation}
que es la Ec.~(12) del artículo. Como $E=K+U$ y $K$ es la energía de un gas ideal
de $f$ grados de libertad, $C_B=n\kB+C$; el calor específico configuracional por
espín que se compara con el canónico es $c=C/N\kB=(C_B-n\kB)/N\kB$.

Note que por la desigualdad de Jensen, $\avg K\avg{1/K}\ge1$, de modo que el lado
derecho de~\eqref{eq:CB-der} es $\le1$ y $C_B\ge n\kB$: el estimador garantiza $C\ge0$.

%---------------------------------------------------------------------
\subsection{Susceptibilidades}
%---------------------------------------------------------------------
\emph{A energía constante.} Con $\avg M=\sum_sM(E-U)^{n-1}\Theta/Z_{n-1}$, el lema
(con $k=n-1$) y $\partial_BM=0$:
\begin{equation}
\chi_E\equiv\Bigl(\frac{\partial\avg M}{\partial B}\Bigr)_E
=(n-1)\frac{\sum_sM^2K^{n-2}\Theta}{Z_{n-1}}-\avg M(n-1)\frac{\sum_sMK^{n-2}\Theta}{Z_{n-1}}
=\Bigl(\frac{dN}{2}-1\Bigr)\Bigl[\Bigl\langle\frac{M^2}{K}\Bigr\rangle-\avg M\Bigl\langle\frac MK\Bigr\rangle\Bigr],
\label{eq:chiE}
\end{equation}
que es la Ec.~(13). Del mismo modo,
\begin{equation}
\Bigl(\frac{\partial\avg M}{\partial E}\Bigr)_B=(n-1)\Bigl[\Bigl\langle\frac MK\Bigr\rangle-\avg M\Bigl\langle\frac1K\Bigr\rangle\Bigr],
\qquad
\Bigl(\frac{\partial\avg K}{\partial B}\Bigr)_E=n\avg M-(n-1)\avg K\Bigl\langle\frac MK\Bigr\rangle .
\label{eq:aux}
\end{equation}
\emph{Relación $\chi_T$--$\chi_E$ en el microcanónico.} Viendo $\avg M$ como función de
$(E,B)$ y luego de $(T,B)$ a través de $E(T,B)$,
\begin{equation}
\chi_T=\chi_E+\Bigl(\frac{\partial\avg M}{\partial E}\Bigr)_B\Bigl(\frac{\partial E}{\partial B}\Bigr)_T,
\qquad
\Bigl(\frac{\partial E}{\partial B}\Bigr)_T=-\frac{(\partial T/\partial B)_E}{(\partial T/\partial E)_B}
=-C_B\Bigl(\frac{\partial T}{\partial B}\Bigr)_E ,
\end{equation}
donde se usó la regla cíclica. Con $T=\avg K/n\kB$ y~\eqref{eq:aux},
$(\partial E/\partial B)_T=\frac{2C_B}{dN\kB}\bigl[(\frac{dN}{2}-1)\avg K\avg{M/K}-\frac{dN}{2}\avg M\bigr]$, y por tanto
\begin{equation}
\chi_T=\chi_E+\Bigl[\frac{dN}{2}-1\Bigr]\Bigl[\Bigl\langle\frac MK\Bigr\rangle-\avg M\Bigl\langle\frac1K\Bigr\rangle\Bigr]
\Bigl[\frac{2C_B}{dN\kB}\Bigr]
\Bigl\{\Bigl[\frac{dN}{2}-1\Bigr]\avg K\Bigl\langle\frac MK\Bigr\rangle-\frac{dN}{2}\avg M\Bigr\},
\end{equation}
que es la Ec.~(14).

\emph{Relación $\chi_E$--$\chi_T$ en el canónico.} Ahora $\avg M$ es función de $(T,B)$:
$\chi_E=\chi_T+(\partial\avg M/\partial T)_B(\partial T/\partial B)_E$, con
$(\partial T/\partial B)_E=-(\partial E/\partial B)_T/C_B$. De~\eqref{eq:E-can},
$(\partial E/\partial B)_T$ es igual a $(\partial\avg U/\partial B)_T$, dada por~\eqref{eq:dUdB-T}. Usando~\eqref{eq:dMdT},
\begin{equation}
\chi_E=\chi_T-\frac{1}{C_B\kB T^2}\bigl(\avg{MU}-\avg M\avg U\bigr)
\Bigl[\frac{1}{\kB T}\bigl(\avg{MU}-\avg M\avg U\bigr)-\avg M\Bigr],
\end{equation}
que es la Ec.~(15).

%---------------------------------------------------------------------
\subsection{Equivalencia de ensambles y origen de la diferencia}
\label{sec:equiv}
%---------------------------------------------------------------------
Sea $g(U)$ la densidad de estados de los espines y $S_c(U)=\kB\ln g(U)$. La
distribución microcanónica de la energía potencial es
$P_E(U)\propto g(U)(E-U)^{n-1}$. Su máximo $U_*$ cumple
$S_c'(U_*)/\kB=(n-1)/(E-U_*)\equiv1/\kB T$ y, como $S_c''=\partial_U(1/T)=-1/(T^2C)$,
\begin{equation}
\frac{\partial^2\ln P_E}{\partial U^2}\Big|_{U_*}=-\frac{1}{\kB T^2C}-\frac{n-1}{(E-U_*)^2}
=-\frac{1}{\kB T^2}\Bigl[\frac1C+\frac{1}{(n-1)\kB}\Bigr].
\end{equation}
En aproximación gaussiana,
\begin{equation}
\avg{\delta U^2}_E\simeq\kB T^2\,\frac{C\,C_K}{C+C_K},\qquad C_K\simeq n\kB=\frac{dN\kB}{2},
\label{eq:dU2}
\end{equation}
es decir, las fluctuaciones de energía están \emph{suprimidas} respecto al valor
canónico $\kB T^2C$: el gas ficticio es un baño de capacidad finita. Pero al llevar
$\avg{\delta K^2}=\avg{\delta U^2}$ a~\eqref{eq:CB-der} con
$\avg K\avg{1/K}\simeq1+\avg{\delta K^2}/\avg K^2$ y $\avg K=n\kB T$:
\begin{equation}
\frac{n\kB}{C_B}\simeq1-(n-1)\frac{C}{n\,(n\kB+C)}=\frac{n\kB}{n\kB+C}+O\Bigl(\frac1N\Bigr)
\quad\Longrightarrow\quad C_B\simeq n\kB+C .
\end{equation}
Por tanto la fórmula~\eqref{eq:CB-der} da el mismo $C$ que el ensamble canónico al
orden dominante: la equivalencia de ensambles se cumple. La diferencia aparece en los
términos de orden siguiente, que involucran las derivadas $\partial C/\partial T$ y
$\partial^2C/\partial T^2$ (términos no gaussianos de $P_E$) y que crecen cerca del
máximo de $C$, donde $C$ varía rápidamente con $T$. Para la semana 15 cuantificaremos
estos términos y su dependencia en $n$ y en $N$, que es la explicación del
16~\% observado.

%=====================================================================
\section{Identificación del ensamble estadístico (preliminar)}
%=====================================================================
\begin{itemize}[leftmargin=1.2em,itemsep=1pt]
\item \textbf{Microestados:} $(\bm p,s)\in\mathbb R^{dN}\times\{-1,+1\}^N$; las
variables de interés son las $2^N$ configuraciones de espín.
\item \textbf{Distribución:} uniforme sobre la hipersuperficie $H(\bm p,s)=E$
(postulado de igual probabilidad a priori); su marginal en $s$ es~\eqref{eq:WE}, que
\emph{no} es uniforme sobre los estados de espín con $U$ fijo.
\item \textbf{Restricciones:} $E$, $N$ y $B$ fijos; $K=E-U\ge0$.
\item \textbf{Equilibrio:} sí; las cadenas de Markov se construyen para que $W_E$ (o
$W_T$) sea su distribución estacionaria, y las propiedades se miden tras la
equilibración.
\item \textbf{Relación con otros ensambles:} $n=1$ da $W_E\propto\Theta(E-U)$, que es
el ensamble del demonio de Creutz con un solo demonio de energía no negativa; $n\to\infty$
con $K/n$ fijo da el canónico [Ec.~\eqref{eq:acc}].
\end{itemize}

%=====================================================================
\section{Código preliminar}
%=====================================================================
El archivo \texttt{codigo/ising\_microcanonico.py} (Python, con \texttt{numba} para
los ciclos sobre espines) implementa:
\begin{enumerate}[itemsep=1pt]
\item \textbf{Metropolis canónico}: barrido secuencial, aceptación
$\min[1,e^{-\Delta U/T}]$, con $\Delta U=2S_ih_i\in\{0,\pm4,\pm8\}$. Estimador
de $c$: Ec.~\eqref{eq:C-can}.
\item \textbf{MC microcanónico}: mismo barrido, aceptación~\eqref{eq:acc} evaluada como
$\exp[(n-1)\ln(K'/K)]$ para evitar desbordes ($n-1=899$); se parte del estado
completamente ordenado ($U=-2N<E$). Estimadores: Ecs.~\eqref{eq:T} y~\eqref{eq:CB-der}.
\item Protocolo del artículo: cinco corridas independientes por punto, con semilla fija;
la incertidumbre es la desviación estándar entre corridas. Las energías
microcanónicas son las de entrada de la Tabla~II.
\end{enumerate}

\paragraph{Prueba de verificación interna (red $4\times4$).}
Las $2^{16}=65\,536$ configuraciones se enumeran para obtener la densidad de estados
exacta $g(U)$. Con ella se calculan, sin aproximación,
$c_{\rm can}(T)$ a partir de~\eqref{eq:C-can} y la curva microcanónica
$\bigl(T(E),c(E)\bigr)$ a partir de $Z_k(E)=\sum_Ug(U)(E-U)^k\Theta(E-U)$,
Ecs.~\eqref{eq:T} y~\eqref{eq:CB-der}. Además se verifica que la solución exacta de
Kaufman para una red finita $m\times n$ con fronteras periódicas~\cite{kaufman49,ferdinand69},
\begin{gather}
Z=\tfrac12(2\sinh2K)^{mn/2}\sum_{i=1}^4Z_i,\qquad K=\beta J,\\
Z_{1,2}=\prod_{r=0}^{n-1}2\left\{\begin{smallmatrix}\cosh\\\sinh\end{smallmatrix}\right\}\Bigl(\tfrac m2\gamma_{2r+1}\Bigr),\quad
Z_{3,4}=\prod_{r=0}^{n-1}2\left\{\begin{smallmatrix}\cosh\\\sinh\end{smallmatrix}\right\}\Bigl(\tfrac m2\gamma_{2r}\Bigr),\nonumber\\
\cosh\gamma_l=\cosh2K\coth2K-\cos\frac{\pi l}{n},\qquad\gamma_0=2K+\ln\tanh K,\nonumber
\end{gather}
reproduce la enumeración (coinciden $c$ y $U$ a $10^{-10}$), de modo que puede
usarse como referencia canónica \emph{exacta} para la red $30\times30$ (curva gris
de la Fig.~\ref{fig:30}).

\begin{figure}[!htbp]
\centering
\figura[width=0.62\textwidth]{figuras/verificacion_4x4.pdf}
\caption{Prueba de verificación: red $4\times4$. Las curvas son exactas (enumeración
de los $2^{16}$ estados); los puntos son MC con $5\times4\times10^5$ barridos.
El MC microcanónico reproduce la curva exacta microcanónica, que difiere
fuertemente de la canónica porque con $N=16$ el baño ficticio es muy pequeño.}
\label{fig:4x4}
\end{figure}

%=====================================================================
\section{Resultados preliminares}
%=====================================================================
Con el mismo número de barridos de medición que el artículo ($5\times10^6$ por punto,
en cinco corridas) pero con una equilibración más corta ($10^5$ barridos), el código
preliminar reproduce las Tablas~I y II (Tabla~\ref{tab:30} y Figs.~\ref{fig:30}--\ref{fig:E}):
\begin{itemize}[leftmargin=1.2em,itemsep=1pt]
\item El calor específico canónico coincide con el del artículo y con el valor exacto
de Kaufman para la red $30\times30$ dentro de una o dos desviaciones estándar en los
nueve puntos. La solución exacta confirma de forma independiente la Tabla~I
(p.\ ej.\ $c_{\rm ex}(2.30)=1.8702$ frente a $1.8712$ del artículo).
\item El calor específico microcanónico en el máximo da $c=2.171\pm0.022$ en
$T=2.2997$, igual al $2.1707$ del artículo: la diferencia entre ensambles del
16~\% se reproduce. La $T$ microcanónica coincide con la Tabla~II en la cuarta
cifra en ocho de los nueve puntos.
\item \textbf{Discrepancia abierta en $E=394.08$:} obtenemos $T=2.1019$ y
$c=0.950\pm0.004$, frente a $T=2.1086$ y $c=0.9713$ en el artículo. Nuestro valor de
$T$ es consistente con la Tabla~I del propio artículo, donde $E=394.09$ corresponde a
$T=2.100$; si fuera $T=2.1086$, con $C_B=N(1+c)\approx1.96\times900$ la energía
tendría que ser unas 15 unidades mayor. Además, en nuestros datos el $c$
microcanónico queda por debajo del canónico a ambos lados del máximo (como en
$T=2.0$ y $2.2$), y el valor del artículo rompe ese patrón. Esto sugiere un error en esa
fila de la Tabla~II, pero antes de afirmarlo hay que descartar una equilibración
insuficiente (cerca de $T=2.1$ la magnetización se invierte sólo cada varios millones
de barridos, según los autores) y estimar mejor las barras de error.
\end{itemize}
Estas barras de error (desviación estándar de cinco corridas) son en sí mismas
inciertas; para la entrega final se complementarán con análisis por bloques y
tiempos de autocorrelación.

\tablaopcional{resultados/tabla_resultados.tex}

\begin{figure}[!htbp]
\centering
\figura[width=\textwidth]{figuras/calor_especifico_30x30.pdf}
\caption{Calor específico configuracional para la red $30\times30$. Símbolos
rellenos claros: valores de las Tablas~I y II del artículo; símbolos huecos: este
trabajo ($5\times10^6$ barridos de medición por punto y $10^5$ de equilibración);
curva gris: solución exacta canónica de Kaufman para $30\times30$.
Izquierda: compárese con la Fig.~2 del artículo; derecha: Fig.~3.}
\label{fig:30}
\end{figure}

\begin{figure}[!htbp]
\centering
\figura[width=0.55\textwidth]{figuras/energia_30x30.pdf}
\caption{Energía total $E$ contra $T$ en ambos ensambles (compárese con la Fig.~1 del
artículo).}
\label{fig:E}
\end{figure}

%=====================================================================
\section{Plan hacia la entrega final (semana 15)}
%=====================================================================
\begin{enumerate}[itemsep=1pt]
\item Convertir el script en un cuaderno Jupyter reproducible de principio a fin.
\item Estudiar la convergencia: autocorrelación integrada de $U$ y $K$ cerca de
$T\approx2.3$, número de barridos de equilibración y estimación del error por bloques,
además de la desviación entre corridas.
\item Barrido fino en $2.20\le T\le2.40$ para reproducir la Fig.~3 completa, y la red
$60\times60$ en $T\approx2.28$.
\item Cuantificar la diferencia entre ensambles con la expansión de la
Sec.~\ref{sec:equiv} y estudiar su dependencia en $n$ (número de grados de libertad
ficticios) y en $N$.
\item Posible extensión: densidad de estados exacta de la red finita~\cite{beale96}
para obtener también el $c$ microcanónico exacto de la red $30\times30$.
\end{enumerate}

%=====================================================================
\appendix
\section*{Anexo: registro preliminar de uso de asistentes de IA}
%=====================================================================
\begin{itemize}[leftmargin=1.2em,itemsep=1pt]
\item \textbf{Qué se pidió:} a Claude (Anthropic) se le pidió un borrador en \LaTeX{}
de este avance y un código preliminar a partir del artículo y del enunciado del
proyecto.
\item \textbf{Qué produjo:} este documento, las derivaciones de la Sec.~3 y el script
\texttt{ising\_microcanonico.py}.
\item \textbf{Cómo se verificó:} el código incluye la comparación con la enumeración
exacta $4\times4$ y con la solución de Kaufman; las derivaciones reproducen las
Ecs.~(9) y (11)--(15) del artículo.
\item \textbf{Pendiente por el grupo:} rehacer a mano cada derivación, revisar
cada línea del código y registrar aquí los errores que se encuentren y cómo se
detectaron.
\end{itemize}

\begin{thebibliography}{9}
\bibitem{ray96} J.~R.~Ray y C.~Frelechoz, ``Microcanonical ensemble Monte Carlo method
for discrete systems'', \textit{Phys.\ Rev.\ E} \textbf{53}, 3402 (1996).
\bibitem{ray91} J.~R.~Ray, \textit{Phys.\ Rev.\ A} \textbf{44}, 4061 (1991).
\bibitem{metropolis53} N.~Metropolis, A.~Rosenbluth, M.~Rosenbluth, A.~Teller y
E.~Teller, \textit{J.\ Chem.\ Phys.} \textbf{21}, 1087 (1953).
\bibitem{creutz83} M.~Creutz, ``Microcanonical Monte Carlo simulation'',
\textit{Phys.\ Rev.\ Lett.} \textbf{50}, 1411 (1983).
\bibitem{lpv67} J.~L.~Lebowitz, J.~K.~Percus y L.~Verlet, \textit{Phys.\ Rev.}
\textbf{153}, 250 (1967).
\bibitem{kaufman49} B.~Kaufman, ``Crystal statistics. II. Partition function evaluated
by spinor analysis'', \textit{Phys.\ Rev.} \textbf{76}, 1232 (1949).
\bibitem{ferdinand69} A.~E.~Ferdinand y M.~E.~Fisher, ``Bounded and inhomogeneous Ising
models. I. Specific-heat anomaly of a finite lattice'', \textit{Phys.\ Rev.}
\textbf{185}, 832 (1969).
\bibitem{beale96} P.~D.~Beale, ``Exact distribution of energies in the two-dimensional
Ising model'', \textit{Phys.\ Rev.\ Lett.} \textbf{76}, 78 (1996).
\end{thebibliography}

\end{document}
